% ---------------------------------------------------------------------------
% ADD THIS SECTION TO sections/appendix.tex
%
% It is the paragraph "A number we withdrew" moved out of Section 2, expanded
% slightly. Section 2 now carries a two-line pointer to it instead, so the
% credibility question is answered in full but is not sitting in the critical
% path of the paper's first result.
%
% Requires \label{app:withdrawn}, which main.tex references twice.
% ---------------------------------------------------------------------------

\section{A number we withdrew}
\label{app:withdrawn}

An earlier internal report put the depth artefact of \S\ref{sec:depth} at
\textbf{46 percentage points} of apparent compositional loss, in the hypothesised
direction, with a depth--maturity correlation of $\rho = -0.92$ and the two arms
$4.3\times$ apart in median depth. Those figures were quoted throughout the
project for weeks and shaped decisions about which analyses to prioritise.

\textbf{We reproduced none of them.} Re-derived from the committed loader, every
sign reverses and the arms come out $0.96\times$ apart rather than $4.3\times$.
The producing script is absent from the repository and the labeller it measured
has since been retired, so the discrepancy resists attribution without guessing
at code nobody kept. We cannot say whether the original was a bug, a different
cell-selection rule, or a transcription error.

\paragraph{What we withdraw and what we keep.} We withdraw the three figures. We
keep the argument, which does not depend on them: it rests on the panel's
construction (\S\ref{sec:intro}, the leakage invariant forcing an
absence-defined label) and on the bound analysis of \S\ref{sec:ceiling}, both of
which are derivable without data. The surviving quantitative claim in
\S\ref{sec:depth} --- $\rho = +0.33$ against a bound of $0.52$, both computed on
the depleted arm --- is re-derived from the committed result table under the
pairing discipline of \S\ref{sec:ceiling}.

\paragraph{Why this is in the paper at all.} Two reasons. First, a reader
comparing this paper against the project's earlier internal circulation would
otherwise find numbers that no longer appear, with no explanation. Second, it is
the cleanest instance in the project of the failure mode the paper is about ---
a quantity that everyone trusted, that no guard was ever pointed at, and that
turned out to be the one number nobody could re-derive. Every other number in
this paper is now read by the figure code directly from a versioned table
carrying a commit hash and a fixed seed, which is a change this episode caused.

\paragraph{Consequence for the reader.} \S\ref{sec:depth}'s empirical claim is
weaker than the withdrawn version and we state it as such: a residual
correlation at 63\% of its attainable bound, on one cohort, after correction. The
structural claims of \S\ref{sec:ceiling}, \S\ref{sec:controls} and
\S\ref{sec:oracle} are unaffected, since none of them is an effect size.
